\uuid{NVxb}
\exo7id{600}
\titre{exo7 600}
\auteur{gourio}
\organisation{exo7}
\datecreate{2001-09-01}
\isIndication{false}
\isCorrection{false}
\chapitre{Suite}
\sousChapitre{Suite récurrente linéaire}
\cpge{mpsi-mp2i.nombres-reels-suites.i, pcsi.nombres-reels-suites.g, ptsi.nombres-reels-suites.g, mpsi-mp2i.nombres-reels-suites.d, pcsi.nombres-reels-suites.c, ptsi.nombres-reels-suites.c, mpsi-mp2i.analyse-asymptotique.c, pcsi.analyse-asymptotique.c, ptsi.analyse-asymptotique.c, bcpst1.suites-reelles, bcpst1.limites-continuite.b}
\module{Analyse}
\niveau{L1}
\difficulte{}

\contenu{
\question{
\'Etudier la suite $(u_{n})_{n\in \Nn}$ d\'{e}finie par $u_{0}=1 $ et
$\forall n\in \Nn$ $u_{n+1}=u_{n}+\frac{2}{u_{n}}.$
En utilisant $v_{n}=\frac{u_{n}^{2}}{4},$ donner un \'{e}quivalent de $%
u_{n}.$
\emph{Indication} : on montrera que $\lim\limits_{n\rightarrow \infty
}v_{n+1}-v_{n}=1, $ on en d\'{e}duira un \'{e}quivalent de $v_{n} $ puis de
$u_{n}.$
}
}
